\documentclass{article}
\usepackage{geometry}[left=1cm,right=1cm]
\usepackage{eso-pic,graphicx} % Required for inserting images
\usepackage[default]{frcursive}

\title{Stormy weather and Magical auroras\\
\large A report of the happening around Embergate's N-E mountains}
\date{\today}

\AddToShipoutPictureBG{\includegraphics[width=\paperwidth,height=\paperheight]{../../Y7.png}}


\begin{document}
    \maketitle

    \section*{Known information}

    "We have been getting some strange reports over the last week or so of thick clouds of various colors on the mountain peak to the north east. Normally this would not concern us, but the terrain has started causing minor quakes and abnormally cold winds for this time of year to come down.

    We need a few adventurers to muscle through the steep trails to see what is going on up there." 

    \section*{Research}

    Arcana pass info: The energy surrounding the tomb appears to have been part of a much larger network of magical seals. The seal beneath the temple was only one point in that system, and the sudden burst suggests that something deeper was disturbed when it was broken. You can also tell that the magenta energy burst was extremely old. Its magical signature does not resemble any common school of magic or known ritual technique. In fact, some of the markings and magical patterns around the tomb appear to have been designed less to contain energy and more to keep something from being reached or awakened.

    Nature pass info: The strange storms were not simply a weather phenomenon. The mountain's natural environment had been reacting to something beneath it. For months, animals had been avoiding the upper slopes. Birds stopped nesting near the peak. Predators moved toward lower elevations, and even hardy mountain plants showed signs of stress despite otherwise favorable conditions. The strange part is that the disturbance seemed to be coming from below the mountain rather than from the summit itself. When the tomb was opened and the magenta energy erupted into the sky, the mountain's natural disturbances suddenly began to subside. The storms cleared almost unnaturally quickly, as though something had been released from the surface tension that had been disturbing the region. The beam could have been fatal to the party, and the region, but the aurora seems to have  appeared in the sky at the right moment to absorb this power in an effort to protect Yaza and their companion.

    \section*{Conclusion}

    From the magical signature the phenomenon from the mountain is outside modern classification of magic, this suggests the origin predate even antique magicks. 

    The preluding storm gathering seems to have been a consequence of gathering energies beneath the raided tomb rather than an intentional magical effect of some kind.

    Similar to static being held within a cumulonimbus cloud the vegetation and fauna surrounding the area noticed and fled where they could and withered where they could not.

    The tomb opening was akin to lighting striking and equalizing the charge. 

    The nature of the barrier which was or held within the tomb seemed to be of a concealing nature rather than a protective layer. 

    This origin beneath the mountain, be it conscious, or a natural phenomenon would have rather than captured, hidden from the outside world. Perhaps to mitigate such an event which occurred the other day. 

    A further investigation of similar positions ought to be funded to understand what happened at the mountain.

    As for the seemingly protective aurora, without an arcane reading it is up to speculation of its nature. Any option between a freak magic surge, divine intervention or the temporary, spontanious acquire of immense arcane ability of the adventurer present.


    
    \end{document}